\section{Data, Provenance and Software Architecture}
\label{app:data}

\subsection{Data sources}

\begin{table}[ht]
\centering
\small
\begin{tabularx}{\linewidth}{>{\raggedright\arraybackslash}p{3.3cm}>{\raggedright\arraybackslash}p{3.4cm}X}
\toprule
Source & Used for & Access and terms \\
\midrule
Cheng et al.\ 2019 supplementary data \cite{cheng2019combinations} & Interactome and drug--target table (default) & Open-access article supplement \todo{confirm reuse terms of the supplementary files} \\
BioGRID \cite{oughtred2021biogrid}, STRING \cite{szklarczyk2023string} & Alternative interactomes (shared with sister project) & Free download; STRING under CC BY 4.0 \todo{confirm BioGRID terms} \\
ChEMBL \cite{zdrazil2024chembl}, DGIdb \cite{freshour2021dgidb} & Supplementary drug--target evidence & ChEMBL under CC BY-SA 3.0; DGIdb free download \todo{confirm DGIdb terms} \\
DrugBank \cite{knox2024drugbank} & Identifier cross-references only (optional) & Academic licence; raw data not redistributed in the repository \\
AD GWAS \cite{bellenguez2022new} & AD module & Published supplementary tables \\
AD pipeline reviews \cite{cummings2025pipeline,cummings2026pipeline} & Curated candidate list & Open-access articles \\
\bottomrule
\end{tabularx}
\caption{Data sources. Raw files are never committed when their licence forbids it; the repository stores download scripts, version strings and SHA-256 checksums instead.}
\end{table}

\begin{plainwords}
Every number the solvers use must be traceable: either computed by a script from a cited source, set as a documented choice in a settings file, or generated randomly for testing (and labeled so). Nobody types in a biological number by hand.
\end{plainwords}

Every numeric value in the objective is one of: \emph{measured} (derived by a documented computation from a cited source), \emph{assumption} (a modeling choice in the configuration, e.g.\ $\lambda$, $z_{\max}$), or \emph{synthetic} (generated for testing, with the generator seed). There are no hand-assigned biological numbers.

\subsection{Schemas}

\paragraph{\texttt{candidates.csv}} (one row per drug; the interface with the sister project)
\begin{lstlisting}
drug_id         stable internal id (e.g., DB00843 or CHEMBL502)
name            preferred name
source_list     curated | sister | scaling_pool
targets         semicolon-separated Entrez Gene IDs present in the interactome
n_targets       integer
d_c, z, p_emp   proximity, z-score, empirical p-value (blank if missing)
n_random        random samples used for z
scorer          "sister@<git hash>" | "reference@<git hash>"
eligible        true/false;  ineligible_reason   free text code (NO_TARGETS, Z_NONNEG, MISSING_Z, NOT_IN_LCC)
provenance      reference to the rows of targets.csv and the scoring run id
\end{lstlisting}

\paragraph{\texttt{targets.csv}} columns \texttt{drug\_id}, \texttt{entrez\_id}, \texttt{source\_db}, \texttt{source\_version}, \texttt{evidence\_type}, \texttt{in\_interactome}.

\paragraph{Instance (problem) files} \texttt{instance.json} plus binary arrays:
\begin{lstlisting}
{ "instance_id": "...", "family": "small|large|synthetic",
  "n": 20, "k": 3, "lambda": 1.0, "A": 0.53,   # penalty weight, used only by QAOA
  "drug_ids": [...],                       # row order = variable order = qubit order
  "files": {"w": "w.npy", "P": "P.npy", "w_int": "w_i32.npy", "P_int": "lamP_i32.npy"},
  "int_scale_bits": 20, "P_layout": "col-major, n_pad=24, 32-byte aligned",
  "sources": {"candidates": "<sha256>", "separation": "<sha256>", "config": "<sha256>"} }
\end{lstlisting}

\paragraph{\texttt{ising.json}} $h$ (dense), $J$ (sparse list of $(i,j,J_{ij})$), offset $C$, scale factor, and the convention string \texttt{"x=(1-z)/2; measured 1 = selected; qubit i = row i of drug\_ids"}.

\paragraph{Run records (\texttt{runs.jsonl})} One JSON object per run: instance id, solver, platform, mode, seed, chains, sweeps or QAOA settings, best set, solver energy, judge energy, valid, found-best, timings by stage (\texttt{t\_prep}, \texttt{t\_build}, \texttt{t\_h2d}, \texttt{t\_kernel}, \texttt{t\_d2h}, \texttt{t\_total}), counters (proposals, accepts, circuit evaluations, shots), peak memory, software versions, git hash.

\paragraph{Files saved along the way} candidate table; shortlist; separation table; penalty table; QUBO and Ising files; temperature schedule and acceptance table; settings files; run records; final combination report (drug names, targets, $z$, pairwise $s$, sources, caveats).

\subsection{Not changing the sister project's scores by accident}

If the sister project supplies scores, its output file is only read, never written, and its SHA-256 hash and configuration are copied into every instance file. A regression test (T12) recomputes $z$ for five drugs with the reference scorer at the same settings and fails if the values differ beyond Monte-Carlo tolerance, which detects silent changes in definitions on either side. No module in this project writes to the candidate scores.

\subsection{Architecture and directory structure}

The core is C++17/CUDA (performance-critical, required by the course) with NASM for the AVX2 kernel; orchestration, data preparation and QAOA are Python (CUDA-Q's primary interface and the sister project's tooling). Solvers are command-line programs that read an instance file and append run records, so every solver, including QAOA, has the same contract.

\begin{lstlisting}
ad-combo/
  data/            raw/ (git-ignored, filled by scripts)  frozen/ (manifests, checksums)
  prep/            01_load_interactome.py  02_targets.py  03_ad_module.py
                   04_reference_scorer.py  05_shortlist.py  06_separation.py
  build/           07_build_instance.py  08_build_ising.py  09_quantize.py
  core/            include/{instance.hpp, record.hpp, philox.hpp, judge.hpp}
                   exact_cpu.cpp  sa_scalar.cpp  sa_avx2.cpp  asm/g_update_avx2.asm
                   intrin/g_update_intrin.cpp  cli/combo_solve.cpp
  cuda/            exact_search.cu  sa_v1.cu  sa_v2.cu
  quantum/         qaoa.py  decode.py  (stretch: qaoa_xy.py  dicke.py)
  bench/           configs/*.yaml  run_small.py  run_large.py  timing.py
  analysis/        figures.py  tables.py
  tests/           test_formulation.py  test_ising.py  test_exact.cpp  test_sa_crossplatform.cpp ...
  docs/            parallel_algorithms.md  results.md  (course deliverables)
\end{lstlisting}

\subsection{Interfaces}

\begin{lstlisting}[language=C++]
struct Instance { int n, k, n_pad; double lambda; double A;
                  std::vector<float> w, P;          // P column-major, padded
                  std::vector<int32_t> w_i, lamP_i; // deterministic mode
                  std::vector<std::string> drug_ids; std::string id, hash; };
struct SolverConfig { std::string solver;   // exact | sa
                      std::string platform; // scalar | avx2 | avx2_intrin | cuda_v1 | cuda_v2
                      Mode mode; uint64_t seed; int chains; int sweeps; double T0, T1; };
struct RunRecord { std::vector<int> best; double e_solver, e_judge; bool valid;
                   Timings t; Counters c; uint64_t seed; };
std::vector<RunRecord> solve(const Instance&, const SolverConfig&);
double judge(const Instance&, const std::vector<int>& selection);   // float64 reference
\end{lstlisting}
The Python side exposes \texttt{solve(instance\_dir, config) -> list[RunRecord]}, calling the CLI for classical solvers and CUDA-Q for QAOA, and writes the same JSON record.

\paragraph{Configuration (YAML)}
\begin{lstlisting}
instance: instances/small_n20_k3_lam1.0
solver: sa
platform: cuda_v1
mode: deterministic
seed: 20261011
chains: 4096
sweeps: 2000
schedule: {type: geometric, T0: auto_80pct, T1: auto_0.1pct}
timing: {warmup: 2, repeats: 10}
\end{lstlisting}

\subsection{Data flow}

\begin{center}
\begin{tikzpicture}[node distance=0.55cm and 0.6cm, box/.style={draw,rounded corners=2pt,align=center,font=\footnotesize,minimum height=0.85cm,text width=2.5cm,fill=gray!5}, gpu/.style={box,fill=green!8}, q/.style={box,fill=blue!6}, >=Stealth]
\node[box] (raw) {Interactome, targets, AD genes, candidate list};
\node[box,right=of raw] (cand) {\texttt{candidates.csv} ($z$, targets)};
\node[box,right=of cand] (sep) {Separation $s$, penalties $P$};
\node[box,right=of sep] (inst) {Instance files, Ising file};
\node[gpu,below=0.9cm of inst] (enum) {Exact search (CPU/GPU)};
\node[gpu,left=of enum] (sa) {SA (scalar, AVX2, CUDA)};
\node[q,left=of sa] (qaoa) {QAOA (CUDA-Q $\to$ cuStateVec)};
\node[box,below=0.9cm of sa,text width=5.5cm] (judge) {Reference evaluator (judge) $\to$ \texttt{runs.jsonl} $\to$ analysis};
\node[box,above=0.6cm of cand, fill=orange!8] (sis) {Sister project output (optional)};
\draw[->] (raw)--(cand); \draw[->] (cand)--(sep); \draw[->] (sep)--(inst); \draw[->,dashed] (sis)--(cand);
\draw[->] (inst)--(enum); \draw[->] (inst.south west)--(sa.north east); \draw[->] (inst.south west)--(qaoa.north east);
\draw[->] (enum.south)|-(judge.east); \draw[->] (sa)--(judge); \draw[->] (qaoa.south)|-(judge.west);
\end{tikzpicture}
\end{center}
